\section{Model Architecture Summary}

In Section \ref{model_architecture} a variety of machine-learning architectures and techniques were discussed. In this section, it will be made explicit the exact machine-learning architectures that will be evaluated. There will be four main categories made up of different variations.

\subsection{SVM with TF-IDF}

\noindent
In this model, features will be generated using the the primitive text-frequency inverse-document-frequency methodology. Features generated by tf-idf will be fed into a linear Support Vector Machine classification head. \\

\noindent 
The sci-kit learn Python library will be used both for the tf-idf and SVM classifier functionality. \\

\noindent
This model was chosen as one of our architectures as it represents arguably the most simplistic way this problem could be approached. In other words, the performance of this model should be treated as the `base-line' performance for which all of the additional more-sophisticated models should at least reach. 

\subsection{Pre-trained SVM-Transformer Hybrid}

\noindent
In this model, features will be extracted using CodeBERT or UniXcoder (both configurations will be tested) in there base pre-trained form (\emph{no} additional fine-tuning). Sentence embeddings will be pooled from the contextualised token embeddings (mean and CLS will be tested). An SVM classifier will then classify on these sentence embeddings. \\

\noindent
To load CodeBERT and UniXcoder, as well as to extract the features and pool the token embeddings, the machine-learning framework PyTorch will be used. As with other models, the linear SVM classifier functitonality will come from sci-kit learn. \\

\noindent
The purpose of this model is two-fold. Firstly, it would be good to know how CodeBERT and UniXcoder perform with no fine-tuning. The results can be compared to the performance of the model when it is fine-tuned, and from that an estimate to how much performance was gained from fine-tuning can be established. Secondly, due to the dataset being quite small, it is possible there are too few training examples to properly train a neural-network classification head. Recall that only the \emph{encoder} has been pre-trained, whereas the classification-head parameters will be trained from scratch on our dataset alone. With this constraint, it is possible that a more simple, linear classification head may perform better (i.e. a SVM classifier).

\subsection{Fine-tuned SVM-Transformer Hybrid}

\noindent
This architecture is identical to the previous pre-trained SVM-transofrmer hybrid, except the encoders will first be fine-tuned on our dataset (using a neural-network classification-head) and then a linear SVM classifier will be trained with features extracted from the training dataset using the fine-tuned encoders. To clairfy: the training dataset for the SVM classifier and encoders are identical. \\

\noindent
As with the similar model using pre-trained encoders, PyTorch will be used the encoders, and sci-kit learn will be used for the SVM classifier.\\

\noindent
This approach will be tested as it is believed fine-tuning the encoder first will enable the encoders to generate better embeddings for the SVM classifier. Furthermore, because the dataset is small, a linear SVM classifier may perform better than a neural-network classifier.\\

\subsection{Fine-tuned Transformer with neural-network classifier}

\noindent
This architecture is most reflective of how CodeBERT and UniXcoder were intended to be used. Pooled sentence embeddings created by the encoders are fed into a neural-network classification head, which generates a probability distribution for the likelihood of the input data being one of th 4 labels. The label with the highest priority is selected.\\

\noindent
PyTorch will be used to implement this model.\\

\noindent
This approach will be tested as it is the most sophisticated of all 4 models that will be tested. It is believed this model will perform the best. 

\subsection{Training Settings}

This section details the important hyperparameters used for the machine-learning models.

\subsubsection{SVM Classifier}

\begin{itemize}
    \item \textbf{Loss function}: hinge loss
    \item \textbf{Penalty:} L2 Regularisation
    \item \textbf{Lagrangian Multiplier:} $e^{-3}$
    \item \textbf{No. Epochs:} $5$
\end{itemize}

\subsubsection{CodeBERT and UniXcoder}

\begin{itemize}
    \item \footnote{This low training batch-size is due to VRAM limitations of the testing machine.}\textbf{Training batch size:} $4$
    \item \textbf{Validation batch size:} $32$
    \item \textbf{Learning Rate: } $5e^{-5}$
    \item \textbf{Max gradient norm: } $1.0$
    \item \footnote{It was experimentally determined that no extra accuracy or decrease in loss was obtained for any additional iterations}\textbf{No. Epochs:} $3$
\end{itemize}

\section{Performance on Validation Dataset}

In this section, the performance of the \emph{trained} models will be evaluated on the validation dataset. As established in Section \ref{dataset_limitations}, these results will be inflated compared to a real-world use-case. \\

\noindent
The various tables below show the performance of each model (and their multiple cofigurations, if applicable) on the validation dataset.

% \subsection{SVM with tf-idf}

\begin{table}[H]
    \begin{center}
        \begin{tabular}{c|c|c}
            \textbf{Feature Generator} & \textbf{Classifier} & \textbf{Accuracy} \\
            \hline
            TF-IDF            & SVM        & 0.66    
        \end{tabular}
    \end{center}
    \caption{Validation performance of SVM-classifier with tf-idf features model}
    \label{tab:validation_perfomance_svm_tfidf}
\end{table}

% \subsection{Pre-trained SVM-Transformer Hybrid}

\begin{table}[H]
    \begin{center}
        \begin{tabular}{c|c|c}
            \textbf{Feature Generator} & \textbf{Pooling Type} & \textbf{Accuracy }\\
            \hline
            CodeBERT & CLS & 0.27 \\
            CodeBERT & Mean & 0.39 \\
            UniXcoder & CLS & 0.78 \\
            UniXcoder & Mean & 0.76 \\
        \end{tabular}
    \end{center}
    \caption{Validation performance of SVM-classifier with tf-idf features model}
    \label{tab:validation_perfomance_pre_svm}
\end{table}

% \subsection{Fine-tuned SVM-Transformer Hybrid}

\begin{table}[H]
    \begin{center}
        \begin{tabular}{c|c|c}
            \textbf{Feature Generator} & \textbf{Pooling Type} & \textbf{Accuracy }\\
            \hline
            CodeBERT & CLS & 0.79 \\
            CodeBERT & Mean & 0.79 \\
            UniXcoder & CLS & 0.79 \\
            UniXcoder & Mean & 0.78 \\
        \end{tabular}
    \end{center}
    \caption{Validation performance of SVM-classifier with tf-idf features model}
    \label{tab:validation_perfomance_fine_svm}
\end{table}

% \subsection{Fine-tuned Transformer with neural-network classifier}

\begin{table}[H]
    \begin{center}
        \begin{tabular}{c|c|c|c}
            \textbf{Feature Generator} & \textbf{Pooling Type} & \textbf{Accuracy} & \textbf{Loss} \\
            \hline
            CodeBERT & CLS & 0.79 & 1.08\\
            CodeBERT & Mean & 0.77 & 1.49\\
            UniXcoder & CLS & 0.80 & 1.08\\
            UniXcoder & Mean & 0.79 & 1.26\\
        \end{tabular}
    \end{center}
    \caption{Validation performance of SVM-classifier with tf-idf features model}
    \label{tab:validation_perfomance_fine_nn}
\end{table}

